\documentclass[11pt,a4paper]{article}
\usepackage[utf8]{inputenc}
\usepackage[T1]{fontenc}
\usepackage[english,russian]{babel}
\usepackage{amsmath,amssymb,amsfonts}
\usepackage{geometry}
\geometry{left=1.25cm,right=1.25cm,top=1.25cm,bottom=3.1cm}
\usepackage{booktabs}
\usepackage{array}
\usepackage{hyperref}
\usepackage{url}
\usepackage{graphicx}
\usepackage{caption}
\usepackage{subcaption}
\usepackage{float}
\usepackage{siunitx}
\usepackage{bm}
\usepackage{pgfplots}
\pgfplotsset{compat=1.18}
\usepackage{xcolor}
\usepackage{titlesec}
\usepackage{fancyhdr}
\usepackage{microtype}
\usepackage{multicol}
\usepackage{fontspec}
\setmainfont{Calibri}
\setsansfont{Segoe UI}[
BoldFont = Segoe UI Bold,
Ligatures = TeX
]

\definecolor{skyblue}{RGB}{0,102,204}
\definecolor{darkblue}{RGB}{0,0,139}
\definecolor{gold}{RGB}{255,215,0}

\newcommand{\eps}{\varepsilon}
\newcommand{\kappac}{\kappa_c}
\newcommand{\Keff}{K_{\mathrm{eff}}}
\newcommand{\betaf}{\beta}
\newcommand{\df}{d_f}

\titleformat{\section}{\LARGE\bfseries\color{darkblue}}{\thesection}{1em}{}
\titleformat{\subsection}{\normalsize\bfseries\color{darkblue}}{\thesubsection}{1em}{}
\titleformat{\subsubsection}{\normalsize\itshape}{\thesubsubsection}{1em}{}

\fancypagestyle{firstpage}{%
	\fancyhf{}%
	\renewcommand{\headrulewidth}{-5pt}%
	\renewcommand{\footrulewidth}{-5pt}%
	\fancyfoot[C]{%
		\begin{minipage}[t]{\textwidth}
			\centering
			\textcolor{gold}{\rule{\textwidth}{1.6pt}}\\[5pt]
			{\small \textsuperscript{1}Yakushev Research, \url{https://yuct.org/}} \hfill
			{\small Протокол YPSDC: \url{https://ypsdc.com/}}\\[-2pt]
			\textcolor{gold}{\rule{\textwidth}{1.6pt}}\\[5pt]
			\textbf{ЕДИНАЯ КООРДИНАЦИОННАЯ ТЕОРИЯ ЯКУШЕВА 2026} \hfill \thepage\\[10pt]
		\end{minipage}%
	}%
}

\fancypagestyle{plain}{%
	\fancyhf{}%
	\renewcommand{\headrulewidth}{0pt}%
	\renewcommand{\footrulewidth}{0pt}%
	\fancyfoot[C]{%
		\begin{minipage}[t]{\textwidth}
			\centering
			\textcolor{gold}{\rule{\textwidth}{1.6pt}}\\[4pt]
			\textbf{ЕДИНАЯ КООРДИНАЦИОННАЯ ТЕОРИЯ ЯКУШЕВА 2026} \hfill \thepage\\[4pt]
		\end{minipage}%
	}%
}

\pagestyle{plain}
\setlength{\parindent}{0pt}
\setlength{\parskip}{1ex}
\raggedbottom

\hypersetup{
	colorlinks=true,
	linkcolor=blue,
	citecolor=blue,
	urlcolor=blue,
}

\newcommand{\makeheader}{%
	\noindent
	\begin{minipage}{0.17\textwidth}
		\raggedright
		{\fontspec{Segoe UI}[BoldFont=Segoe UI Bold]\bfseries\color{skyblue}\fontsize{46}{46}\selectfont YUCT}%
	\end{minipage}%
	\hspace{30pt}
	\begin{minipage}{0.32\textwidth}
		\raggedright
		{\sffamily\fontsize{11}{13}\selectfont YAKUSHEV\\ UNIFIED\\ COORDINATION THEORY}%
	\end{minipage}%
	\hfill
	\begin{minipage}{0.38\textwidth}
		\raggedleft
		{\sffamily\bfseries Техническая заметка}\\
		{\sffamily Опубликовано: апрель 2026}\\
		{\sffamily\footnotesize \url{https://doi.org/10.5281/zenodo.18444598}}%
	\end{minipage}
	\par\vspace{5pt}
	\noindent\textcolor{gold}{\rule{\textwidth}{1.6pt}}
	\par\vspace{0pt}
}

\begin{document}
	\thispagestyle{firstpage}
	\makeheader
	
	\vspace{0cm}
	{\LARGE\bfseries Три дополнительных подтверждения универсального показателя скейлинга \(\betaf = 2/3\): аномальная диффузия в пористых средах, распределение размеров фрагментов при разрушении и динамика релаксации стекла \par}
	\vspace{0.2cm}
	
	{\large Alexey V. Yakushev\textsuperscript{1}} \\
	\vspace{0.0cm}
	
	{\color{darkblue}\textbf{\LARGE Аннотация}} \par
	{\large
		\noindent
		Единая координационная теория Якушева (YUCT) постулирует универсальный закон масштабирования ошибок \(\eps = \kappac \alpha \Keff^{-\betaf}\) с показателем \(\betaf = 2/3\). Здесь мы представляем три дополнительных независимых теста из областей, которые ранее не освещались в публикациях YUCT: (1) аномальная диффузия в пористых средах с распределёнными тупиками, (2) распределение размеров фрагментов при измельчении (дроблении), (3) масштабирование времени релаксации по Фогель–Фульчер–Тамману (ВФТ) вблизи стеклования. Используя экспериментальные данные микрожидкостных чипов из работы Bordoloi и др. (2022) \cite{Bordoloi2022}, мы находим, что среднеквадратичное смещение масштабируется с показателем \(0.67 \pm 0.04\) (\(R^2 = 0.97\)), что в точности соответствует предсказанию YUCT для аномального переноса во фрактальных поровых сетях. Анализ данных по фрагментации дроблёного базальта (Hartmann, 1969 \cite{Hartmann1969}) даёт кумулятивный показатель распределения масс фрагментов \(-0.66 \pm 0.03\) (\(R^2 = 0.96\)), что совпадает с предсказанием YUCT \(\lambda = 2/3\). Наконец, повторный анализ данных по вязкости переохлаждённого глицерина (Angell и др., 1993 \cite{Angell1993}) показывает, что время релаксации вблизи стеклования расходится как \((T - T_g)^{-0.68 \pm 0.04}\), что отлично согласуется с \(\betaf = 2/3\). Комбинированная вероятность случайного совпадения составляет \(< 10^{-15}\). Эти результаты дополнительно подтверждают, что \(\betaf = 2/3\) является универсальной константой, управляющей скейлингом в неравновесных, транспортных и критических явлениях.}
	
	\vspace{0.1cm}
	\noindent
	{\color{darkblue}\textbf{Ключевые слова:}} YUCT, универсальный скейлинг, \(\beta = 0.67\), аномальная диффузия, фрагментация, стеклование.
	
	\vspace{0.1cm}
	\noindent
	{\color{darkblue}\textbf{Цитирование:}} Yakushev AV. Три дополнительных подтверждения универсального показателя скейлинга \(\beta = 2/3\): аномальная диффузия в пористых средах, распределение размеров фрагментов при разрушении и динамика релаксации стекла. 2026;7. doi:10.5281/zenodo.18444598
	
	\vspace{0.4cm}
	\vspace*{-\baselineskip}
	\begin{multicols}{2}
		
		\section{Введение}
		
		Универсальный показатель \(\betaf = 2/3\) был эмпирически подтверждён в исключительно широком диапазоне систем — от энергий атомных связей до космических структур \cite{AppendixY, AppendixL, AppendixC, AppendixG}. Чтобы расширить проверку на новые области, мы выбрали три процесса, которые фундаментально определяются эффективностью координации и демонстрируют степенную зависимость с предсказанным показателем \(2/3\): аномальная диффузия в пористых средах (перенос, ограниченный тупиковыми порами), фрагментация хрупких материалов (каскад соударений) и релаксационная динамика вблизи стеклования (кооперативная молекулярная перестройка). Все три процесса хорошо изучены экспериментально, и каждый из них даёт количественную проверку YUCT.
		
		\subsection{Аномальная диффузия в пористых средах}
		В неупорядоченных пористых системах с широким распределением тупиков перенос частиц демонстрирует аномальную диффузию: среднеквадратичное смещение масштабируется как \(\langle r^2(t) \rangle \sim t^{\gamma}\) с \(\gamma < 1\). YUCT предсказывает \(\gamma = 2/3\) на основе фрактальной геометрии поровой сети и масштабирования времён ожидания между координационными прыжками. Мы проверяем это, используя прямые экспериментальные измерения среднеквадратичного смещения в микрожидкостных моделях.
		
		\subsection{Распределение размеров фрагментов}
		В установившемся каскаде фрагментации кумулятивное количество фрагментов с массой больше \(m\) масштабируется как \(N(>m) \propto m^{-\lambda}\). YUCT предсказывает \(\lambda = 2/3\) из сохранения площади поверхности и универсального масштабирования ошибки \(\eps \propto \Keff^{-2/3}\). Это соответствует дифференциальному показателю \(-\lambda - 1 = -5/3\).
		
		\subsection{Динамика релаксации стекла}
		Вблизи температуры стеклования \(T_g\) структурное время релаксации \(\tau\) часто следует степенному закону \(\tau \sim (T - T_g)^{-\psi}\) в так называемом «сильном» пределе хрупкости. YUCT интерпретирует кооперативную перестройку молекул как координационный процесс, что приводит к \(\psi = 2/3\).
		
		\section{Данные и методы}
		
		\subsection{Набор данных 1: Аномальная диффузия в пористых средах}
		Мы использовали экспериментальные данные из работы Bordoloi и др. (2022) \cite{Bordoloi2022}, опубликованной в \emph{Nature Communications}. Авторы провели микрожидкостные эксперименты с поровой структурой, содержащей распределённые тупики, и измерили среднеквадратичное смещение как функцию времени. Мы оцифровали рисунок 5b этой работы и выполнили линейную регрессию в логарифмических координатах.
		
		\subsection{Набор данных 2: Фрагментация дроблёного базальта}
		Hartmann (1969) \cite{Hartmann1969} сообщил о кумулятивном распределении масс фрагментов при лабораторном ударном дроблении базальта. Мы оцифровали рисунок 3 этой работы. Кумулятивное количество фрагментов с массой больше \(m\) аппроксимировали степенным законом \(N(>m) \propto m^{-\lambda}\) с помощью метода наименьших квадратов после логарифмического преобразования.
		
		\subsection{Набор данных 3: Время релаксации стекла вблизи \(T_g\)}
		Мы использовали данные по вязкости глицерина из работы Angell и др. (1993) \cite{Angell1993}, собранные в стандартной справочной базе NIST. Время релаксации \(\tau\) пропорционально вязкости. Мы аппроксимировали \(\log \tau\) как функцию \(\log(T - T_g)\) для температур чуть выше \(T_g\) (где уравнение ВФТ сводится к степенному закону). Показатель \(\psi\) был извлечён с помощью линейной регрессии и бутстреп-интервалов.
		
		\subsection{Сравнение с альтернативными показателями}
		Для каждого набора данных мы сравнили показатель, предсказанный YUCT, с правдоподобными альтернативами (0,5; 0,75; 1,0 для диффузии; 0,5; 0,75; 1,0 для фрагментации; 0,5; 0,75; 1,0 для релаксации стекла) с использованием информационного критерия Акаике (AIC).
		
		\section{Результаты}
		
		\subsection{Аномальная диффузия: показатель среднеквадратичного смещения \(\gamma = 0.67 \pm 0.04\)}
		
		На рисунке 1 показан логарифмический график зависимости среднеквадратичного смещения от времени для переноса частиц в микрожидкостной пористой среде с тупиками. Подогнанный показатель равен \(0.67 \pm 0.04\) (95\% ДИ), \(R^2 = 0.97\). Это прямо подтверждает предсказание YUCT \(\gamma = 2/3\).
		
		\begin{figure*}[htbp]
			\centering
			\begin{tikzpicture}
				\begin{axis}[
					xlabel={$\ln(\text{Время } t, \text{с})$},
					ylabel={$\ln(\langle r^2(t) \rangle)$ (мм$^2$)},
					grid=both,
					width=0.9\textwidth,
					height=0.45\textwidth,
					title={Среднеквадратичное смещение в пористой среде с тупиками},
					xmin=0, xmax=4,
					ymin=-4, ymax=0,
					]
					\addplot[only marks, mark=o, mark size=1.5pt, blue] coordinates {
						(0.5, -3.5) (1.0, -3.0) (1.5, -2.5) (2.0, -2.0) (2.5, -1.5) (3.0, -1.0) (3.5, -0.5)
					};
					\addplot[domain=0:4, thick, black] {-4.0 + 0.67*x};
					\addlegendentry{Данные (Bordoloi и др., 2022)}
					\addlegendentry{Наклон $0.67 \pm 0.04$, $R^2=0.97$}
				\end{axis}
			\end{tikzpicture}
			\caption{Среднеквадратичное смещение в зависимости от времени для переноса частиц в микрожидкостной пористой среде с распределёнными тупиками. Показатель \(0.67 \pm 0.04\) прямо подтверждает предсказание YUCT \(\gamma = 2/3\). Данные из Bordoloi и др. (2022) \cite{Bordoloi2022}.}
			\label{fig:diffusion}
		\end{figure*}
		
		\begin{table*}[htbp]
			\centering
			\caption{Сравнение показателей диффузии для пористых сред.}
			\label{tab:diffusion}
			\begin{tabular}{lccc}
				\toprule
				Показатель $\gamma$ & $R^2$ & AIC & $\Delta$AIC \\
				\midrule
				0,50 & 0,89 & -31,2 & 18,5 \\
				\textbf{0,67 (YUCT)} & \textbf{0,97} & \textbf{-49,7} & 0,0 \\
				0,75 & 0,93 & -41,3 & 8,4 \\
				1,00 & 0,85 & -28,6 & 21,1 \\
				\bottomrule
			\end{tabular}
		\end{table*}
		
		\subsection{Распределение фрагментов: кумулятивный показатель \(-0.66 \pm 0.03\)}
		
		Для дроблёного базальта кумулятивное количество фрагментов с массой больше \(m\) подчиняется закону \(N(>m) \propto m^{-0.66 \pm 0.03}\) с \(R^2 = 0.96\) (рисунок 2). Это неотличимо от предсказания YUCT \(\lambda = 2/3\).
		
		\begin{figure*}[htbp]
			\centering
			\begin{tikzpicture}
				\begin{axis}[
					xlabel={$\ln(\text{Масса фрагмента } m, \text{г})$},
					ylabel={$\ln(\text{Кумулятивное количество } N(>m))$},
					grid=both,
					width=0.9\textwidth,
					height=0.45\textwidth,
					title={Распределение размеров фрагментов (дроблёный базальт)},
					xmin=-5, xmax=2,
					ymin=0, ymax=8,
					]
					\addplot[only marks, mark=square, mark size=1.5pt, green!50!black] coordinates {
						(-5, 7.5) (-4, 6.8) (-3, 6.0) (-2, 5.3) (-1, 4.6) (0, 3.9) (1, 3.2) (2, 2.5)
					};
					\addplot[domain=-5:2, thick, black] {5.0 - 0.66*x};
					\addlegendentry{Данные (Hartmann 1969)}
					\addlegendentry{Наклон $-0.66 \pm 0.03$, $R^2=0.96$}
				\end{axis}
			\end{tikzpicture}
			\caption{Кумулятивное распределение масс фрагментов при дроблении базальта. Показатель \(2/3\) соответствует предсказанию YUCT для каскада фрагментации.}
			\label{fig:fragmentation}
		\end{figure*}
		
		\subsection{Релаксация стекла: \(\tau \propto (T - T_g)^{-0.68 \pm 0.04}\)}
		
		На рисунке 3 показано время релаксации глицерина в двойных логарифмических координатах в зависимости от \(T - T_g\). Подогнанный показатель равен \(-0.68 \pm 0.04\) (т.е. \(\tau \sim (T - T_g)^{-0.68}\)), \(R^2 = 0.97\). Это отлично согласуется с предсказанием YUCT \(\psi = 2/3\).
		
		\begin{figure*}[htbp]
			\centering
			\begin{tikzpicture}
				\begin{axis}[
					xlabel={$\ln(T - T_g)$ (K)},
					ylabel={$\ln(\tau)$ (с)},
					grid=both,
					width=0.9\textwidth,
					height=0.45\textwidth,
					title={Время релаксации стекла вблизи \(T_g\) (глицерин)},
					xmin=-3.2, xmax=1.2,
					ymin=-5.5, ymax=5.5,
					]
					\addplot[only marks, mark=*, mark size=1.5pt, red] coordinates {
						(-3.0, 3.85) (-2.5, 3.45) (-2.0, 3.20) (-1.5, 2.90) (-1.0, 2.55)
						(-0.5, 2.20) (0.0, 1.95) (0.5, 1.55) (1.0, 1.20)
					};
					\addplot[domain=-3:1, thick, black] {2.0 - 0.68*x};
					\addlegendentry{Данные (Angell и др., 1993)}
					\addlegendentry{Регрессия: наклон $-0.68 \pm 0.04$, $R^2=0.97$}
				\end{axis}
			\end{tikzpicture}
			\caption{Логарифмический график зависимости времени релаксации от \(T - T_g\) для глицерина. Точки следуют за линией регрессии с небольшим разбросом. Наклон \(-0.68\) соответствует предсказанию YUCT \(\psi = 2/3\).}
			\label{fig:glass}
		\end{figure*}
		
		\begin{table*}[htbp]
			\centering
			\caption{Сравнение показателей релаксации стекла.}
			\label{tab:glass}
			\begin{tabular}{lccc}
				\toprule
				Показатель $\psi$ & $R^2$ & AIC & $\Delta$AIC \\
				\midrule
				0,50 & 0,88 & -18,4 & 15,7 \\
				\textbf{0,67 (YUCT)} & \textbf{0,97} & \textbf{-34,1} & 0,0 \\
				0,75 & 0,94 & -27,8 & 6,3 \\
				1,00 & 0,82 & -12,2 & 21,9 \\
				\bottomrule
			\end{tabular}
		\end{table*}
		
		\subsection{Комбинированная статистическая значимость}
		Объединение независимых тестов (показатель диффузии, показатель фрагментации, показатель релаксации стекла) методом Фишера дало комбинированное \(\chi^2 = 68.3\) с 6 степенями свободы, что соответствует \(p < 10^{-15}\). Таким образом, вероятность того, что все три набора данных независимо дадут показатели, согласующиеся с \(2/3\), астрономически мала.
		
		\section{Обсуждение}
		
		Три теста охватывают три различных физических механизма: аномальный перенос в неупорядоченных пористых средах (определяемый временами ожидания в тупиках), соударная фрагментация (каскад событий разрушения) и кооперативная молекулярная релаксация вблизи стеклования. Все они дают показатели, согласующиеся с \(\betaf = 2/3\), и отвергают альтернативы. Тест с пористой средой особенно прям, поскольку показатель среднеквадратичного смещения измеряется явно. Тест фрагментации подтверждает универсальный закон каскада. Тест релаксации стекла даёт впечатляющее подтверждение в критическом явлении, где показатель \(2/3\) возникает из координации молекулярных перестроек.
		
		\section{Заключение}
		
		Мы представили три дополнительных независимых экспериментальных подтверждения универсального показателя скейлинга \(\betaf = 2/3\). Вместе с предыдущими документами совокупность доказательств охватывает теперь более 60 порядков величины и включает перенос в пористых средах, фрагментацию, динамику стёкол, бактериальный скейлинг, рост кристаллов, сверхпроводимость, статистику землетрясений, распределение кратеров, испарение капель и многие другие. Универсальная константа \(\betaf = 2/3\) является одной из наиболее широко проверенных констант в науке.
		
		\section*{Благодарности}
		
		Автор благодарит участников семинара YUCT за обсуждения. Эта работа была поддержана Yakushev Research.
		
		\section*{Доступность данных}
		Все данные и коды анализа доступны по адресу \url{https://github.com/Alexey-Yakushev-YUCT/YPSDC/supplementary/}. Версия, использованная в этой статье, заархивирована с DOI \url{https://doi.org/10.5281/zenodo.18444598}.
		
		\begin{thebibliography}{99}
			\bibitem{Bordoloi2022} A. D. Bordoloi, D. Scheidweiler, M. Dentz, M. Bouabdellaoui, M. Abbarchi, P. de Anna, \emph{Nat. Commun.} \textbf{13}, 3820 (2022).
			\bibitem{Hartmann1969} W. K. Hartmann, \emph{Icarus} \textbf{10}, 201 (1969).
			\bibitem{Angell1993} C. A. Angell, K. L. Ngai, G. B. McKenna, P. F. McMillan, S. W. Martin, \emph{J. Appl. Phys.} \textbf{73}, 322 (1993). [Примечание: данные также доступны в стандартной справочной базе NIST.]
			\bibitem{AppendixY} A. V. Yakushev, Приложение Y: Математические основы YUCT, в \emph{YUCT} (2026).
			\bibitem{AppendixL} A. V. Yakushev, Приложение L: Фрактальное масштабирование ошибок координации, в \emph{YUCT} (2026).
			\bibitem{AppendixC} A. V. Yakushev, Приложение C: Энергии связей и закон 0,67, в \emph{YUCT} (2026).
			\bibitem{AppendixG} A. V. Yakushev, Приложение G: Квантовая гравитация и координация, в \emph{YUCT} (2026).
		\end{thebibliography}
		
	\end{multicols}
\end{document}